\documentclass{article}
\usepackage[T1]{fontenc}
\usepackage{lmodern}
\usepackage{graphicx}
\usepackage{amsmath,amssymb}
\usepackage[a4paper,margin=2cm]{geometry}
\usepackage[absolute,overlay]{textpos}
\setlength{\TPHorizModule}{1mm}
\setlength{\TPVertModule}{1mm}
\usepackage[indonesian]{babel}
\usepackage{hyperref}
\setlength{\emergencystretch}{2em}
% unit: Halaman 13

\begin{document}

% === Halaman 13 (llm) ===

Bukti bahwa pesan $m$ dapat diungkap kembali dari pasangan cipherteks $a$ dan $b$:

Dari persamaan (2): $a = g^k \pmod p \rightarrow g^k \equiv a \pmod p$

Pangkatkan kedua ruas dengan $x$ : $g^{xk} \equiv a^x \pmod p$

Dari persamaan (3): $b = y^k m \pmod p \rightarrow y^k m \equiv b \pmod p$

Dari persamaan (1): $y = g^x \pmod p \rightarrow g^x \equiv y \pmod p$,

maka

\begin{align*}
 b/a^x &\equiv y^k m / a^x \pmod p \\
 &\equiv g^{xk} m / g^{xk} \pmod p \\
 &\equiv m \pmod p
\end{align*}

yang berarti bahwa plainteks $m$ dapat diungkap kembali dari cipherteks $a$ dan $b$.

\end{document}
